\documentclass[12pt,a4paper]{article}
\usepackage[utf8]{inputenc}
\usepackage[T2A]{fontenc}
\usepackage[russian]{babel}
\usepackage{graphicx}
\usepackage{booktabs}
\usepackage{longtable}
\usepackage{geometry}
\usepackage{float}
\usepackage{fontspec}
\setmainfont{DejaVu Serif}
\setmonofont{DejaVu Sans Mono}
\usepackage{paracol}
\usepackage{src/title}
\usepackage{hyperref}

\usepackage{listings}
\usepackage{xcolor}

\geometry{top=2cm, left=2cm, right=2cm, bottom=2cm}

\lstset{
  basicstyle=\ttfamily\small,
  breaklines=true,
  language=Prolog,
  showstringspaces=false,
}

\title{Рекомендательная система выбора танка на основе базы знаний Prolog\\
\large Лабораторная работа №2. Система поддержки принятия решений}
\author{Ануфриев А.С., группа 465029}
\supervisor{Научный руководитель кафедры программной инженерии и вычислительной техники}
\date{27.10.2026}

\abstract{
Разработана рекомендательная система (система поддержки принятия решений) для выбора боевого танка на основе базы знаний, представленной на языке Prolog. Пользователь вводит запрос на русском языке из трёх частей: что надо вернуть (список / рекомендация / что есть в наличии), объект запроса (танк, уровень, завод, нация) и набор фильтрующих условий (нация, тип техники, уровень, высокая альфа, максимальная скорость, минимальная перезарядка, конкретный танк). Условия накладываются на выборку в порядке их следования во входном строке. Программа на Python парсит строку запроса и посылает логические запросы к базе знаний на SWI-Prolog, которая выступает в роли источника данных (БД). На основе результатов выборки система формирует рекомендацию с обоснованием.
}
\keywords{рекомендательная система, база знаний, Prolog, SWI-Prolog, системы поддержки принятия решений, танки}

\begin{document}
\maketitle

\section*{Введение}
\addcontentsline{toc}{section}{Введение}

Целью данной работы является разработка программы --- рекомендательной системы, которая использует базу знаний для предоставления рекомендаций на основе введённых пользователем данных. В качестве subject-домена выбраны боевые танки Второй мировой войны (база знаний из лабораторной работы №1, файл \texttt{tanks.pl}).

В отличие от первого варианта реализации (чистый Prolog-диалог), в этой работе архитектура разделена: Python-приложение отвечает за парсинг пользовательского ввода и формирование ответа, а база знаний на Prolog (\texttt{tanks.pl} вместе с уровнем доступа \texttt{tank\_db.pl}) служит источником данных, к которому Python посылает логические запросы. Такой подход демонстрирует использование базы знаний как базы данных и логического вывода для фильтрации и ранжирования объектов.

\section*{Анализ требований}
\addcontentsline{toc}{section}{Анализ требований}

К системе предъявлены следующие требования:
\begin{itemize}
\item пользователь вводит запрос одной строкой на русском языке фиксированного формата из трёх частей;
\item строка парсится на факты (режим вывода, объект запроса, условия);
\item выполняются логические запросы к базе знаний на Prolog;
\item система выдаёт рекомендацию (или список) с обоснованием.
\end{itemize}

\subsection*{Формат входного запроса}

Запрос состоит из трёх частей в строгом порядке:

\begin{center}
\begin{tabular}{lll}
\toprule
Часть & Назначение & Варианты \\
\midrule
1 & Что вернуть (режим) & \texttt{выведи} --- список; \texttt{подскажи} --- рекомендация с обоснованием; \texttt{какой есть} --- что есть в наличии \\
2 & Объект & \texttt{танк} (по умолчанию), \texttt{уровень}, \texttt{завод}, \texttt{нация} \\
3 & Условия (3.1--3.7) & нация, тип, уровень, высокая альфа, быстрый, скорострельный, конкретный танк \\
\bottomrule
\end{tabular}
\end{center}

Условия (часть 3) можно объединять; они накладываются на выборку в порядке их следования во входной строке:

\begin{center}
\begin{tabular}{lll}
\toprule
№ & Условие & Допустимые значения \\
\midrule
3.1 & Нация & \texttt{советский}, \texttt{немецкий}, \texttt{американский} \\
3.2 & Тип техники & \texttt{тт} (тяжёлый), \texttt{ст} (средний), \texttt{лт} (лёгкий) \\
3.3 & Уровень & \texttt{7лвл}, \texttt{8 лвл} (также \texttt{лвл 7/8}) \\
3.4 & Высокая альфа & урон больше всех на уровне (по умолчанию 8 уровень) \\
3.5 & Быстрый & скорость больше всех среди текущей выборки \\
3.6 & Скорострельный & перезарядка меньше всех среди текущей выборки \\
3.7 & Конкретный танк & отображаемое имя танка (например, \texttt{ИС-3}) \\
\bottomrule
\end{tabular}
\end{center}

Режимы вывода:
\begin{itemize}
\item \textbf{выведи} --- список всех объектов, прошедших фильтры, с характеристиками;
\item \textbf{подскажи} --- рекомендация одного объекта из прошедших фильтров с обоснованием (по каким условиям он подобран);
\item \textbf{какой есть} --- перечень доступных значений объекта (нации, уровни, заводы или имена танков) среди прошедших фильтров.
\end{itemize}

\section*{База знаний}
\addcontentsline{toc}{section}{База знаний}

База знаний \texttt{tanks.pl} содержит 18 танков 7 и 8 уровней трёх наций (СССР, Германия, США) трёх классов (тяжёлый, средний, ПТ-САУ). Каждый танк описан фактами:
\begin{itemize}
\item \texttt{tank/1} --- идентификатор танка;
\item \texttt{tank\_stats/11} --- уровень, нация, класс, HP, урон/выстрел (альфа), пробитие, броня, имя, калибр, перезарядка;
\item \texttt{built\_by/3} --- производитель и год начала разработки;
\item \texttt{unlocked\_from/2} --- базовый танк предыдущего уровня;
\item \texttt{team/2}, \texttt{team\_title/2} --- комплекты по нациям и уровням.
\end{itemize}

Реализованы правила \texttt{durable/1} (живучесть, HP $\ge$ 1800), \texttt{hard\_hitter/1} (пробитие $\ge$ 220 мм), \texttt{all\_rounder/1}, \texttt{glass\_cannon/1}, \texttt{opponents/2}, \texttt{best\_gun/2}.

Для поддержки условий ``быстрый'' в уровень доступа \texttt{tank\_db.pl} добавлены факты \texttt{speed/2} (ориентировочная максимальная скорость, км/ч) для всех 18 танков.

\section*{Архитектура и реализация}
\addcontentsline{toc}{section}{Архитектура и реализация}

\begin{center}
\begin{tabular}{ll}
\toprule
Компонент & Роль \\
\midrule
\texttt{recommend.py} & Парсер русских запросов, формирование ответов, диалоговый CLI \\
SWI-Prolog (\texttt{swipl}) & Механизм логического вызова, источник данных \\
\texttt{tank\_db.pl} & Уровень доступа к данным поверх \texttt{tanks.pl} \\
\texttt{tanks.pl} & База знаний (факты и правила) \\
\bottomrule
\end{tabular}
\end{center}

Python посылает запрос Prolog через стандартный ввод \texttt{swipl} (ключ \texttt{-t} в используемой сборке SWI-Prolog 10 некорректно обрабатывает цели с переменными, поэтому цель подаётся через stdin), парсит построчный вывод и формирует ответ на русском языке.

Основные предикаты \texttt{tank\_db.pl}, вызываемые из Python:
\begin{center}
\begin{tabular}{ll}
\toprule
Предикат & Назначение \\
\midrule
\texttt{all\_tanks(-Ts)} & все танки базы \\
\texttt{filter\_tanks(+Nation,+Class,+Level,-Ts)} & фильтрация по нации/классу/уровню (\texttt{any} = без ограничения) \\
\texttt{alpha\_max\_tanks(+Level,-Ts)} & танки с максимальным уроном/выстрелом на уровне \\
\texttt{fastest\_tanks(+Ids,-Ts)} & из списка --- с максимальной скоростью \\
\texttt{fastest\_reloading\_tanks(+Ids,-Ts)} & из списка --- с минимальной перезарядкой \\
\texttt{tank\_by\_name(+Name,-Tank)} & поиск по отображаемому имени \\
\texttt{tank\_row(...)} & полная строка характеристик танка \\
\texttt{print\_tank(+Tank)} & строка для парсинга Python, поля через ``|'' \\
\texttt{all\_nations/1}, \texttt{all\_levels/1}, \texttt{all\_manufacturers/1} & справочники \\
\bottomrule
\end{tabular}
\end{center}

Алгоритм обработки запроса:
\begin{enumerate}
\item Загружается вся таблица характеристиков из Prolog (\texttt{print\_tank} для каждого танка).
\item Строка парсится слева направо: сначала режим вывода и объект, затем условия в порядке вхождения.
\item Начальная выборка --- все танки; каждое условие --- отдельный логический запрос к Prolog с пересечением с текущей выборкой.
\item Сверхъятивные условия (\texttt{альфа}, \texttt{быстрый}, \texttt{скорострельный}) применяются к текущей выборке: ``урон больше всех на уровне'' считается на указанном (или 8-м) уровне среди текущих кандидатов; ``быстрый'' и ``скорострельный'' --- по текущей выборке.
\item Результат проецируется по объекту (танки / нации / уровни / заводы) и выводится согласно режиму.
\end{enumerate}

\section*{Примеры запросов и результатов}
\addcontentsline{toc}{section}{Примеры запросов и результатов}

\begin{center}
\begin{tabular}{llp{7cm}}
\toprule
Запрос & Режим & Результат \\
\midrule
\texttt{выведи танк нация советский лвл 8 с высокой альфой} & список & ИСУ-152 (640 урона/выстрел, максимум на 8 уровне среди советских 8-го уровня) \\
\texttt{подскажи танк немецкий лвл 7 быстрый} & рекомендация & Panther I --- максимальная скорость 55 км/ч среди немецких 7-го уровня \\
\texttt{выведи танк американский скорострельный} & список & T20 --- перезарядка 7 с, минимум среди американских танков \\
\texttt{какой есть нации} & наличие & Советский Союз, Германия, США \\
\texttt{выведи танк конкретный танк ИС-3} & список & ИС-3: HP 1800, урон 400, пробитие 225 мм, скорость 37 км/ч \\
\texttt{выведи танк тип лт} & список & не найдено (в базе нет лёгких танков) \\
\texttt{подскажи танк лвл 8 с высокой альфой быстрый} & рекомендация & ИСУ-152 --- высочайшая альфа (640) и лучшая скорость в выборке \\
\bottomrule
\end{tabular}
\end{center}

\section*{Тестирование}
\addcontentsline{toc}{section}{Тестирование}

Проведено 11 сквозных сценариев диалога (входной запрос $\rightarrow$ ответ системы). Все сценарии прошли: рекомендации соответствуют фактам из \texttt{tanks.pl}, пустые результаты обрабатываются сообщением о неудаче, команда выхода завершает программу с кодом 0.

Пример работы (\texttt{подскажи танк немецкий лвл 7 быстрый}):
\begin{lstlisting}
Рекомендую: Panther I
Обоснование: максимальная скорость (55 км/ч).
Panther I (Германия, 7 лвл, средний): HP 1300, урон/выстрел 280, пробитие 198 мм,
броня 100 мм, калибр 88 мм, перезарядка 5.5 с, скорость 55 км/ч
\end{lstlisting}

Пример работы (\texttt{выведи танк тип ст лвл 8}):
\begin{lstlisting}
Найдено 3 танк(ов):
1. Т-44 (Советский Союз, 8 лвл, средний): ... скорость 53 км/ч
2. Panther II (Германия, 8 лвл, средний): ... перезарядка 5.14 с, скорость 46 км/ч
3. Pershing (США, 8 лвл, средний): ... скорость 40 км/ч
\end{lstlisting}

\section*{Заключение}
\addcontentsline{toc}{section}{Заключение}

Реализована рекомендательная система, в которой база знаний на Prolog используется как источник данных, а Python отвечает за парсинг естественноязычного (русского) запроса и представление результатов. Система корректно обрабатывает комбинируемые фильтры в порядке их следования, три режима вывода и способы обоснования рекомендации. Возможные направления развития: расширение базы знаний новыми танками и характеристиками, добавление ранжирования по весу условий, графический интерфейс.

\section*{Приложение А. Листинг tank\_db.pl}
\addcontentsline{toc}{section}{Приложение А. Листинг tank\_db.pl}

\begin{lstlisting}
:- consult('tanks.pl').

speed(is_3, 37). speed(t_44, 53). speed(isu_152, 26).
speed(tiger_i, 45). speed(panther_i, 55). speed(jpanther, 46).
speed(tiger_ii, 34). speed(panther_ii, 46). speed(ferdinand, 22).
speed(t29, 35). speed(t_20, 40). speed(t_25_2, 40).
speed(t32, 35). speed(pershing, 40). speed(t28_prot, 13).
speed(is, 37). speed(t_43, 42). speed(su_152, 28).

filter_tanks(Nation, Class, Level, Ts) :-
    findall(T,
            (tank_stats(T, Lvl, Nat, Cls, _, _, _, _, _, _, _),
             (Nation = any; Nat = Nation),
             (Class = any; Cls = Class),
             (Level = any; Lvl = Level)),
            Ts).

alpha_max_tanks(Level, Ts) :-
    findall(Alpha-Tank,
            tank_stats(Tank, Level, _, _, _, Alpha, _, _, _, _, _),
            Pairs),
    (Pairs = [] -> Ts = [];
     keysort(Pairs, Sorted),
     last(Sorted, MaxAlpha-_),
     findall(Tank, (member(Alpha-Tank, Pairs), Alpha = MaxAlpha), Ts)).

fastest_tanks(Ids, Ts) :-
    findall(Speed-Tank, (member(Tank, Ids), speed(Tank, Speed)), Pairs),
    (Pairs = [] -> Ts = [];
     keysort(Pairs, Sorted),
     last(Sorted, MaxSpeed-_),
     findall(Tank, (member(Speed-Tank, Pairs), Speed = MaxSpeed), Ts)).
\end{lstlisting}

\section*{Приложение Б. Листинг recommend.py (фрагмент парсера)}
\addcontentsline{toc}{section}{Приложение Б. Листинг recommend.py (фрагмент)}

\begin{lstlisting}[language=Python]
def parse_request(text: str) -> dict:
    tokens = tokenize(text)
    mode, entity = None, None
    conditions = []  # (kind, value) в порядке вхождения
    i, n = 0, len(tokens)
    while i < n:
        t = tokens[i]
        if mode is None and t in MODE_ALIASES:
            mode = MODE_ALIASES[t]
        elif mode is None and t == "какой" and i + 1 < n and tokens[i + 1] == "есть":
            mode = "what_exists"; i += 1
        elif entity is None and t in ENTITY_ALIASES:
            entity = ENTITY_ALIASES[t]
        elif t == "с" and i + 2 < n and tokens[i + 1] == "высокой" and tokens[i + 2] == "альфой":
            conditions.append(("alpha", None)); i += 2
        elif t in NATION_ALIASES:
            conditions.append(("nation", NATION_ALIASES[t]))
        ...
        i += 1
    return {"mode": mode or "show", "entity": entity or "tank",
            "conditions": conditions}
\end{lstlisting}

\end{document}
